\begin{apartats}

\begin{tria}{definicio-tasca}
\itemtria{original}{1}{1,25}
Aplicant la definició de derivada, calcula la funció derivada de
\[
f(x)=2x^2-3x .
\]

\begin{solucio}
\[
f'(x)=\lim_{h\to0}\frac{2(x+h)^2-3(x+h)-\left(2x^2-3x\right)}{h}
=\lim_{h\to0}\frac{4xh+2h^2-3h}{h}=\lim_{h\to0}(4x+2h-3)=4x-3 .
\]
\end{solucio}

\itemtria{reconeix-limit}{1}{1,25}
Cadascun dels límits següents és la derivada d'una funció en un punt. Digues de quina funció
i en quin punt, i calcula'l:
\begin{graella}{2}
  \sa \lim_{h\to0}\frac{(2+h)^3-8}{h} & \sa \lim_{h\to0}\frac{e^{h}-1}{h}
\end{graella}

\begin{solucio}
La definició és $f'(a)=\lim_{h\to0}\dfrac{f(a+h)-f(a)}{h}$.\\
i) És $f'(2)$ amb $f(x)=x^3$, perquè $2^3=8$. Com que $f'(x)=3x^2$, el límit val $f'(2)=12$.
També es pot desenvolupar: $\dfrac{12h+6h^2+h^3}{h}=12+6h+h^2\to12$.\\
ii) És $f'(0)$ amb $f(x)=e^{x}$, perquè $e^0=1$. Com que $f'(x)=e^{x}$, el límit val $f'(0)=1$.
\end{solucio}
\end{tria}

\apartat[1,25]{0,75}
Aplicant la definició de derivada, calcula la funció derivada de
\[
g(x)=\sqrt{x+3} .
\]

\begin{solucio}
Multipliquem i dividim pel conjugat:
\[
\begin{aligned}
g'(x)&=\lim_{h\to0}\frac{\sqrt{x+h+3}-\sqrt{x+3}}{h}
=\lim_{h\to0}\frac{(x+h+3)-(x+3)}{h\left(\sqrt{x+h+3}+\sqrt{x+3}\right)}\\
&=\lim_{h\to0}\frac{1}{\sqrt{x+h+3}+\sqrt{x+3}}=\frac{1}{2\sqrt{x+3}} ,
\end{aligned}
\]
per a $x>-3$.
\end{solucio}

\begin{nomesllarg}
\apartat{0,75}
Aplicant la definició de derivada, calcula la funció derivada de
\[
k(x)=\frac{4}{x} .
\]

\begin{solucio}
\[
k'(x)=\lim_{h\to0}\frac{\frac{4}{x+h}-\frac{4}{x}}{h}
=\lim_{h\to0}\frac{4x-4(x+h)}{h\,x(x+h)}
=\lim_{h\to0}\frac{-4}{x(x+h)}=-\frac{4}{x^2} ,
\]
per a $x\neq0$.
\end{solucio}
\end{nomesllarg}

\end{apartats}
